\section{Post-Training}\label{sec:posttraining}

Post-training proceeds in two supervised stages (Table~\ref{tab:stages});
neither uses reinforcement learning or a thinking mode. Standard
supervised fine-tuning (SFT) turns \mossbase{} into \mossinstruct{}, an
offline instruction follower. Realtime-SFT then continues from
\mossinstruct{} (Table~\ref{tab:lineage}) and installs the real-time
interaction paradigm: deciding at every frame whether to speak, staying
silent while nothing needs saying, and revising an answer when the scene
overturns it. Every real-time-specific design choice in \mossvl{} lives
in this final stage, which accounts for under 3\% of the total training
tokens.

\subsection{Supervised Fine-Tuning}\label{sec:sft}

We fine-tune \mossbase{} on 7.6M instruction samples (102.8B tokens)
with the standard next-token cross-entropy loss over assistant
responses, with sequences up to 128K tokens (Table~\ref{tab:stages}).
The samples combine data collected and reorganized from existing
corpora with data synthesized in house, and all of it passes filtering,
deduplication, decontamination against our evaluation suites, and
quality screening before entering the mixture. The
mixture covers general question answering over single images,
multi-image sets, videos, and plain text; perception-centric tasks
including OCR, document understanding, and spatial and temporal
grounding; image and video captioning; and reasoning-centric tasks
spanning multimodal reasoning, mathematics and other academic
disciplines, code, and knowledge-intensive QA, together with identity
data.

\subsection{Realtime-SFT: Learning When to Speak}\label{sec:realtime-sft}

Realtime-SFT teaches the model to treat incoming video as a stream of
decisions rather than a finished artifact. Training samples interleave
text with frames in arrival order: every frame is followed by a decision
slot, and each slot takes one of three forms---\statetok{silence} (keep
watching), \statetok{response} followed by text (speak now), or a reply
that ends with \statetok{silence} (finish speaking). A reply is spread
over consecutive slots frame by frame, emulating rate-limited real-time
output, and a single user turn may contain several separate emissions.
Supporting this costs exactly two new vocabulary entries---the two state
tokens, initialized from the embeddings of semantically related existing
tokens. The speak-or-wait decision itself is ordinary next-token
prediction: whenever the most probable next token is \statetok{silence},
the model waits for the next frame; otherwise it decodes a reply. No
dedicated decision head is attached.

\paragraph{Data characteristics.}
What sets the Realtime-SFT corpus (0.56M samples, ${\approx}$34.8B
tokens; Table~\ref{tab:stages}) apart is that every sample casts the
model in an explicit interaction role rather than a plain QA role:
standing instructions that must fire exactly once when their condition
is met; resident questions whose answers must update as evidence
accumulates; continuous real-time commentary; counting that accumulates
across a stream; probes of whether this is the right moment to speak;
and video-independent dialogue that maintains identity consistency. A
further share of offline QA and general multimodal data preserves
offline ability.

\paragraph{Data construction.}
The corpus draws on two sources. We first collect open-source datasets
for streaming video understanding and subject them to strict filtering
and re-annotation. More important, however, is the data we synthesize
ourselves, targeting the behaviors that existing datasets provide
scarcely or not at all: staying silent until evidence appears, revising
an answer as the scene evolves, and recovering when a new event
interrupts a reply midway. Synthesis follows the caption-driven pipeline
introduced in MOSS-Video-Preview \citep{Wang2026moss}: hierarchical,
densely time-anchored captions are mined for state transitions of a
focal object; each transition yields a question, an immediately
answerable reply, and a trajectory of updates; replies are anchored to
visual moments, laid out over frames with silence in between, and
filtered for quality. This round upgrades three of the pipeline's four
stages. Temporal anchoring is now verified frame by frame against the
actual footage: each reply is assigned the moment its evidence becomes
visible and the moment it stops being valid. Hand-offs are more
natural: a reply overtaken by a new event is rewritten as a
plain-language self-correction instead of being marked with a dedicated
interrupt token. Finally, a quality gate checks every sample against
its frames and keeps only those whose replies are grounded in what is
visible at emission time.

\paragraph{Corpus statistics.}
The interaction-first emphasis is visible in the numbers. Across the
corpus the model is supervised on 2.2M emission decisions, 58.7\% of
which are self-timed rather than prompted by a fresh user question; in
5.1\% of samples the target event never occurs, and the correct
behavior is to stay silent throughout. Streams run at 1\,fps for up to
768 frames (${\approx}$12.8 minutes) per window. The corpus is
decontaminated against our evaluation suites: benchmark videos are
excluded via held-out lists.

\subsection{Mode Control and Training Objective}\label{sec:mode-control}

Mode control in \mossvl{} amounts to a single system prompt. The
streaming and real-time modes share one prompt---real-time operation is
a special case of streaming---and offline inference uses none. One set
of weights thus operates in three inference modes with zero architecture
change. The prompt is reproduced below; the full dialogue template,
including frame and timestamp interleaving, is given in
Appendix~\ref{app:template}.

\begin{promptbox}{The shared streaming / real-time system prompt}
You are a helpful AI assistant specializing in real-time video analysis.
The video streams to you frame by frame. At every frame, you decide
independently whether to respond or stay silent --- output
\statetok{silence} when nothing relevant has happened, and respond when
the visual content warrants it.
\end{promptbox}

Supervision covers only what the assistant controls: reply text and the
two state tokens. System and user turns and the expanded visual tokens
are excluded from the loss. The central difficulty is imbalance: silence
slots vastly outnumber emission decisions, and under uniform weights the
model simply learns to stay silent. We therefore reweight the two state
tokens with a focal factor and inverse-frequency class coefficients,
computing the class statistics over the global batch at every step,
which keeps the class coefficients identical across data-parallel
ranks:
\begin{equation}\label{eq:realtime-loss}
\mathcal{L}=\frac{\sum_i m_i\, w_i\, \ell_i}{\sum_i m_i},\qquad
w_i=\begin{cases}
\alpha_{y_i}\,(1-p_i)^{\gamma} & \text{$y_i$ is a state token,}\\
1 & \text{otherwise,}
\end{cases}
\end{equation}
where $\ell_i$ is the token-level cross-entropy, $m_i$ the supervision
mask, $p_i$ the predicted probability of the target token, and
$\gamma=2$. The coefficient $\alpha_k=(n_s+n_r)/(2\,n_k)$, with $n_s$ and
$n_r$ counting silence and response targets in the current global batch,
equalizes the nominal, pre-focal weight of the two classes; reply text
keeps unit weight. Before focal modulation, the decision to speak and
the decision to stay silent thus carry equal aggregate weight in the
state-token loss, no matter how rare speaking is.

One further masking choice matters specifically in streams. In offline
chat, the token closing an assistant turn marks the end of an exchange;
in a stream, a turn boundary usually means the user interjected while the
world---and the conversation---continue. We therefore also exclude the
assistant's turn-final end token from supervision, so the model never
learns to wrap up merely because a new user turn appears. In a controlled
single-variable comparison, this masking raised emission frequency by
39\% and mean reply length by 68\%.

The behaviors installed here are evaluated quantitatively on four
streaming benchmarks in \S\ref{sec:evaluation} and qualitatively in
Figure~\ref{fig:showcase}.
